\section{Almost reflective dependently typed language:
    \texorpdfstring{$\twoLTT$}{[[2]]LTT}}
\label{twoLTT}

Now we are ready to take on the main challenge: design
a minimal dependently typed reflective language, and
provide internal models of itself to yield universe
polymorphism, external parametricity and do synthetic
mathematics inside it. As a starting point, we take
Martin-L\"of type theory with $\Pi-$,$\Sigma-$,
$\mathrm{Id}-$ types, natural number type and infinite
non-cumulative hierarchy of universes. As discussed in
Section~\ref{no-go-section}, the best practical
solution to reinterpretation here would be to forbid
reductions under it and reinterpret terms as-is, and
constrain models using $\mathrm{Id}$-types. Still, in
comparison to non-dependently typed case, there are
infinitely more problems to solve if we want to define
a truly reflective system:

\begin{enumerate}
    \item If we na\"{\i}vely adopt an approach to
        typing of models from Section~\ref{SSTLC}, in
        order to specify a model, we would need to
        specify countably many models, one per each
        universe. And even if we did, to which universe
        does the type of models belong?
    \item Now that we have computations in types, we
        have to deal with conversion rule under
        reinterpretation:
        \[
            \dfrac{t:T\quad T\equiv T'}{t:T'}
            \stackrel{?}{\leadsto}
            \dfrac{\llbracket t\rrbracket:\llbracket T
            \rrbracket\quad\llbracket T\rrbracket\equiv
            \llbracket T'\rrbracket}{\llbracket t
            \rrbracket:\llbracket T'\rrbracket}
        \]
        The problem is that, after interpretation,
        $\llbracket T \rrbracket$ is only
        \textit{propositionally} equal to $\llbracket
        T'\rrbracket$.
    \item If we want a truly reflective system, we need
        to know what to do with
        \textit{reinterpretation of reinterpretation}:
        \[
            \lambda x.\llbracket\llbracket x
            \rrbracket_1\rrbracket_2
            \quad\text{and}\quad
            \llbracket\lambda x.\llbracket x
            \rrbracket_1\rrbracket_2
        \]
        In the first case, we perform two subsequent
        interpretations on a yet unknown $x$; in the
        second case, we interpret the process of
        interpretation itself. What should this mean?
\end{enumerate}

To deal with universe sizing issues, consider the
following observation: in a system without level
binders, each term uses only a finite number of levels.
Hence to reinterpret this term, we need to only provide
a finite number of universe models, which can together
fit into some larger universe of a finite level. No
universe polymorphism required! However, we now have to
track the maximal level which appears in a term in our
typing judgements.

To interpret conversion rule, we start by first
pinpointing the locations where conversion rule is
actually used; in order to do this ergonomically, we
use bidirectional type checking. At the moment of
conversion, interpretation has enough information to
know both $T$ and $T'$; thus we can translate a chain
of reductions performed during definitional equality
check into a chain of propositional equalities taken
from coherence conditions proven in the definition of a
model. Then, a point of conversion is wrapped in
equality elimination with synthesized witness as a
scrutinee.

In reinterpreting reinterpretation, we actually only
deal with ``double reinterpretation'': cases like
$\lambda x.\llbracket\llbracket x\rrbracket_1
\rrbracket_2$ could be dealt with by \textit{merging}
two models into a new one, a technique reminiscent of
\textit{Collapsing towers of interpreters}
\cite{Amin2017}; it is not required in our system since
the innermost reinterpretation is allowed to compute in
this case. As for interpreting interpretation itself as
in $\llbracket\lambda x.\llbracket x \rrbracket_1
\rrbracket_2$, we explicitly forbid such constructions
by introducing a family of modalities $\Box_\kappa$
meaning ``no interpretation was used in construction of
this term, and all levels appearing in it are not above
$\kappa$''. Terms have to have a boxed type to be
reinterpreted. In the notation for modality and general
typing discipline for it, we largely follow Gratzer et
al \cite{Gratzer2019} with a few important differences
discussed later. We also use $\Box$ to relax a ban on
computations inside reinterpretation by allowing
``splicing'' boxed expressions into other boxed
expressions; inside splices, computation is allowed to
bring terms to the boxed form for further
interpretation. Such a typing discipline brings us very
close to two-level type theories [CITE], but we found
notation via modalities to be more convenient.

\begin{figure}
\centering
\begin{bnf}
$\Gamma$ : Context ::=
    | $\cdot$ : empty context
    | $\Gamma,x\,\colon T$ : variable extension
    | $\Gamma.\,{}_\ell\text{\faLock}$
    : context locking
    ;;
$T$, $S$ : Type ::=
    | $\mathcal{U}_\ell$ : universe
    | $\mathrm{El}(t)$ : universal type family
    | $\mathrm{Nat}_\ell$ : natural number type
    | $(x\,\colon S)\times T$ : $\Sigma$-type
    | $(x\,\colon S)\to T$ : $\Pi$-type
    | $t_1=_T t_2$ : identity type
    | $\Box_\ell\,T$ : type of interpretable terms
    | $t\,\llbracket T\rrbracket$
        : interpretation of types
    ;;
$\ell$, $\kappa$, $j$ : Level :in: $\mathbb{N}$
    : universe levels ;;
$t$, $s$ : Term ::=
    | $x$ // $y$ // $z$ // $\ldots$
        : variables
    | $\mathtt{U}_\ell$ : code of a universe
    | $\mathtt{Nat}_\ell$
        : code of a natural number type
    | $0$ : zero
    | $\mathrm{S}\,t$ : successor
    | $\mathrm{natind}_{x.T}\,t\,\{s_0,s_S\}$
        : induction
    | $\mathtt{Sigma}(x\,\colon s)t$
        : code of a $\Sigma$-type
    | $(s, t)$ : pairing
    | $\pi_1\,t$ // $\pi_2\,t$ : projections
    | $\mathtt{Pi}(x\,\colon s)t$
        : code of a $\Pi$-type
    | $\lambda x.\,t$ : abstraction
    | $t\,s$ : application
    | $\mathtt{Id}\,t\,t_1\,t_2$
        : code of an identity type
    | $\mathrm{refl}$ : reflexivity
    | $\mathrm{J}_{xy.S}\,t\,s$
        : equality elimination
    | $\mathtt{Box}_\ell\,t$ : code of a box type
    | $[t]_\text{\faLock}$ : locking
    | $[t]_\text{\faUnlock}$ : unlocking
    | $t\,\llbracket s\rrbracket$
        : interpretation of terms
    | $t\,\colon T$ : ascription
    ;;
\end{bnf}
\caption{$\twoLTT$ grammar}
\label{twoltt-grammar}
\end{figure}

Grammar rules given in Figure~\ref{twoltt-grammar}
reflect discussion above. Note that there is no
separate type of models $\mathcal{M}$: instead of there
being a single builtin type, there is a family of
model types $\mathcal{M}_{\ell_0,\ldots,\ell_\kappa}^{
    \ell'_0,\ldots,\ell'_\kappa}$ where each concrete
type is \emph{definable} in $\twoLTT$; types in a
family are indexed by a bunch of universe levels whose
exact role will be explained later; a general idea is
that $\ell_j$ is a level which types in $\mathcal{U}_j$
would have after interpretation in such a model.

Typing rules are also a straightforward adaptation of a
general recipe for bidirectional typing of
dependently-typed systems \cite{Thiago2025} and are
given in Figures~\ref{twoltt-wf}-\ref{twoltt-typing}.
Important departure from standard bidirectional
recipe ``constructors are checked, eliminators are
inferred'' is that levels of types are always inferred
since we do not have cumulativity of universes and do
not have a lifting operator, thus we cannot check
levels of $\Sigma$ and $\Pi$. Also note that exact
model type is always inferred, both in interpretation
of types and in interpretation of terms: even though
$\mathcal{M}$ would be some huge nested $\Sigma$-type,
it is still more correct to \textit{match} on it
instead of asking to provide levels explicitly.

As already said, reinterpretation is constrained to act
on boxed types; in addition, we require a non-locked
context for reinterpretation. Together these
constraints guarantee that, indeed, boxed terms do not
contain reinterpretation subterms. An important
departure from $\mathrm{MLTT}_\text{\faLock}$
\cite{Gratzer2019}, aside from adding a family of
modalities, lies in the way context is locked in
$[\_]_\text{\faLock}$ and unlocked in
$[\_]_\text{\faUnlock}$: while locking, we require the
outer context to be unlocked; while unlocking, instead
of erasing all locks, we erase the context up to the
first lock. Effectively, this enforces a constraint
that typing context never contains more than a single
lock. Thus, for example, our $\Box_\ell$ is not a
comonad, but axiom K from modal logic still
holds\footnote{Interestingly enough, a similar
K-modality is first described by Birkedal and others in
``Modal dependent type theory and dependent right
adjoints'' \cite{Birkedal2020}. Coincidentally, they
refine typing judgements with universe levels, too!
They have a single modality, though.}:
\[
    \dfrac{\cdot\vdash_\ell A:\ell_A
    \quad\cdot\vdash_\ell B:\ell_B}{
        \cdot\vdash_\kappa\lambda f x.[[f
        ]_\text{\faUnlock}[x]_\text{\faUnlock}
        ]_\text{\faLock}:\Box_\ell(A\to B)\to
            \Box_\ell A\to\Box_\ell B
    }(\ell_A\sqcup\ell_B\leqslant\ell)
\]
In this way, unlocking truly gains the semantics of
``splicing'' in the sense that it can only contain
boxed values from the ``outer world''.

\begin{figure}
\centering
\begin{mathpar}
    \inferrule[WF-Nil]{}{\cdot\vdash_\ell}
    \and
    \inferrule[WF-Var]{
        \Gamma\vdash_\ell\\
        \Gamma\vdash_\ell T:\kappa
    }{\Gamma,x:T\vdash_\ell}(\kappa\leqslant\ell)
    \and
    \inferrule[WF-Lock]{
        \Gamma\vdash_\ell\\
        \text{\faLock}\not\in\Gamma
    }{\Gamma.\,{}_\ell\text{\faLock}\vdash_\kappa}
    \and
    \inferrule[Level-El]
        {\Gamma\vdash_\ell t\Rightarrow
            \mathcal{U}_\kappa}
        {\Gamma\vdash_\ell\mathrm{El}(t):\kappa}
        (\kappa+1\leqslant\ell)
    \and
    \inferrule[Level-Uni]{\Gamma\vdash_\ell}
        {\Gamma\vdash_\ell\mathcal{U}_\kappa:\kappa+1}
    \and
    \inferrule[Level-Nat]{\Gamma\vdash_\ell}
        {\Gamma\vdash_\ell\mathrm{Nat}_\kappa:\kappa}
    \and
    \inferrule[Level-Sigma]{
        \Gamma\vdash_\ell S:j \\
        \Gamma,x:S\vdash_\ell T:\kappa
    }{\Gamma\vdash_\ell(x:S)\times T:j\sqcup\kappa}
    \and
    \inferrule[Level-Pi]{
        \Gamma\vdash_\ell S:j \\
        \Gamma,x:S\vdash_\ell T:\kappa
    }{\Gamma\vdash_\ell(x:S)\to T:j\sqcup\kappa}
    \and
    \inferrule[Level-Id]{
        \Gamma\vdash_\ell T:\kappa \\
        \Gamma\vdash_\ell t_1\Leftarrow T \\
        \Gamma\vdash_\ell t_2\Leftarrow T
    }{\Gamma\vdash_\ell t_1=_T t_2:\kappa}
    \and
    \inferrule[Level-Box]{
        \text{\faLock}\not\in\Gamma\\
        \Gamma.\,{}_\ell\text{\faLock}\vdash_\kappa T:j
    }{\Gamma\vdash_\ell\Box_\kappa\,T:\kappa}
    (j\leqslant\kappa)
    \and
    \inferrule[Level-TyInterp]{
        \Gamma\vdash_\ell t\Rightarrow\mathcal{M}_{
            \ell_0,\ldots,\ell_\kappa}^{
            \ell'_0,\ldots,\ell'_\kappa}\\
        \text{\faLock}\not\in\Gamma\\
        \Gamma.{}_\ell\text{\faLock}\vdash_\kappa T:j
    }{\Gamma\vdash_\ell t\llbracket T\rrbracket:\ell_j}
    (j\leqslant\kappa)
\end{mathpar}
\caption{
    Context well-formedness \& level placement rules}
\label{twoltt-wf}
\end{figure}

\begin{figure}
\centering
\begin{mathpar}
    \inferrule[Infer-Uni]{\Gamma\vdash_\ell}{
        \Gamma\vdash_\ell\mathtt{U}_\kappa\Rightarrow
            \mathcal{U}_{\kappa+1}
    }
    \and
    \inferrule[Infer-Nat]{\Gamma\vdash_\ell}{
        \Gamma\vdash_\ell\mathtt{Nat}_\kappa\Rightarrow
            \mathcal{U}_\kappa}
    \and
    \inferrule[Check-Zero]{\Gamma\vdash_\ell}{
        \Gamma\vdash_\ell 0\Leftarrow
            \mathrm{Nat}_\kappa
    }
    \and
    \inferrule[Check-Succ]
        {\Gamma\vdash_\ell t\Leftarrow
            \mathrm{Nat}_\kappa}
        {\Gamma\vdash_\ell\mathrm{S}\,t\Leftarrow
            \mathrm{Nat}_\kappa}
    \and
    \inferrule[Infer-NatInd]{
        \Gamma\vdash_\ell t\Rightarrow
            \mathrm{Nat}_\kappa\\
        \Gamma,x:\mathrm{Nat}_\kappa\vdash_\ell T:j\\
        \Gamma\vdash_\ell s_0\Leftarrow
            T[(0:\mathrm{Nat}_\kappa)/x]\\
        \Gamma\vdash_\ell s_S\Leftarrow
                (y:\mathrm{Nat}_\kappa)\to
                (z:T[y/x])\to
                T[(\mathrm{S}\,y:\mathrm{Nat}_\kappa)
                /x]
    }{\Gamma\vdash_\ell
        \mathrm{natind}_{x.T}\,t\,\{s_0,s_S\}
        \Rightarrow T[t/x]
    }(j\leqslant\ell)
    \and
    \inferrule[Infer-Sigma]{
        \Gamma\vdash_\ell s\Rightarrow
            \mathcal{U}_j\\
        \Gamma,x:\mathrm{El}(s)\vdash_\ell t
            \Rightarrow\mathcal{U}_\kappa
    }{\Gamma\vdash_\ell\mathtt{Sigma}(x:s)t\Rightarrow
        \mathcal{U}_{j\sqcup\kappa}}
    \and
    \inferrule[Check-Pair]{
        \Gamma\vdash_\ell s\Leftarrow S\\
        \Gamma\vdash_\ell t\Leftarrow T[(s:S)/x]
    }{\Gamma\vdash_\ell(s,t)\Leftarrow(x:S)\times T}
    \and
    \inferrule[Infer-$\mathrm{Proj}_1$]
        {\Gamma\vdash_\ell t\Rightarrow(x:S)\times T}
        {\Gamma\vdash_\ell \pi_1\,t\Rightarrow S}
    \and
    \inferrule[Infer-$\mathrm{Proj}_2$]
        {\Gamma\vdash_\ell t\Rightarrow(x:S)\times T}
        {\Gamma\vdash_\ell \pi_2\,t\Rightarrow
            T[(\pi_1\,t)/x]}
    \and
    \inferrule[Infer-Pi]{
        \Gamma\vdash_\ell s\Rightarrow
            \mathcal{U}_j\\
        \Gamma,x:\mathrm{El}(s)\vdash_\ell t
            \Rightarrow\mathcal{U}_\kappa
    }{\Gamma\vdash_\ell\mathtt{Pi}(x:s)t\Rightarrow
        \mathcal{U}_{j\sqcup\kappa}}
    \and
    \inferrule[Check-Abs]
        {\Gamma,x:S\vdash_\ell t\Leftarrow T}
        {\Gamma\vdash_\ell\lambda x.\,t\Leftarrow
            (x:S)\to T}
    \and
    \inferrule[Infer-App]{
        \Gamma\vdash_\ell t\Rightarrow (x:S)\to T\\
        \Gamma\vdash_\ell s\Leftarrow S
    }{\Gamma\vdash_\ell t\,s\Rightarrow
        T[(s:S)/x]}
    \and
    \inferrule[Infer-Id]{
        \Gamma\vdash_\ell t\Rightarrow
            \mathcal{U}_\kappa\\
        \Gamma\vdash_\ell t_1\Leftarrow\mathrm{El}(t)
        \\\Gamma\vdash_\ell t_2\Leftarrow\mathrm{El}(t)
    }{\Gamma\vdash_\ell\mathtt{Id}\,t\,t_1\,t_2
        \Rightarrow\mathcal{U}_\kappa}
    \and
    \inferrule[Check-Refl]{t_1\equiv t_2}
        {\Gamma\vdash_\ell\mathrm{refl}\Leftarrow
            t_1=_T t_2}
    \and
    \inferrule[Infer-J]{
        \Gamma\vdash_\ell t\Rightarrow t_1=_T t_2\\
        \Gamma,x:T,y:t_1=_T x\vdash_\ell S:\kappa\\\\
        \Gamma\vdash_\ell s\Leftarrow
            S[(t_1:T)/x,
                (\mathrm{refl}:t_1=_T t_1)/y]
    }{\Gamma\vdash_\ell\mathrm{J}_{xy.S}\,t\,s
        \Rightarrow S[(t_2:T)/x, t/y]
    }(\kappa\leqslant\ell)
    \and
    \inferrule[Infer-Var]{
        \text{\faLock}\not\in\Delta\\
        \Gamma,x:T,\Delta\vdash_\ell
    }{\Gamma,x:T,\Delta\vdash_\ell x\Rightarrow T}
    \and
    \inferrule[Infer-Box]{
        \text{\faLock}\not\in\Gamma\\
        \Gamma.\,{}_\ell\text{\faLock}\vdash_\kappa t
            \Rightarrow\mathcal{U}_j
    }{\Gamma\vdash_\ell\mathtt{Box}_\kappa\,t
        \Rightarrow\mathcal{U}_\kappa}
    (j\leqslant\kappa)
    \and
    \inferrule[Check-Lock]{
        \text{\faLock}\not\in\Gamma\\
        \Gamma.\,{}_\ell\text{\faLock}\vdash_\kappa t
            \Leftarrow T
    }{\Gamma\vdash_\ell[t]_\text{\faLock}\Leftarrow
        \Box_\kappa T}
    \and
    \inferrule[Infer-Unlock]{
        \text{\faLock}\not\in\Delta\\
        \Gamma\vdash_\kappa t\Rightarrow\Box_j T
    }{\Gamma.{}_\kappa\text{\faLock},\Delta
        \vdash_\ell[t]_\text{\faUnlock}\Rightarrow T
    }(j\leqslant\ell)
    \and
    \inferrule[Infer-Interp]{
        \text{\faLock}\not\in\Gamma\\
        \Gamma\vdash_\ell s\Rightarrow\mathcal{M}_{
            \ell_0,\ldots,\ell_k}^{
            \ell'_0,\ldots,\ell'_k}\\
        \Gamma\vdash_\ell t\Rightarrow \Box_\kappa T
    }{\Gamma\vdash_\ell s \llbracket t\rrbracket
    \Rightarrow s \llbracket T\rrbracket
    }
    \and
    \inferrule[Infer-Ascr]{
        \Gamma\vdash_\ell T:\kappa\\
        \Gamma\vdash_\ell t\Leftarrow T
    }{\Gamma\vdash_\ell t:T\Rightarrow T}
    (\kappa\leqslant\ell)
    \and
    \inferrule[Check-Inferred]{
        \Gamma\vdash_\ell t\Rightarrow S\\ S\equiv T
    }{\Gamma\vdash_\ell t\Leftarrow T}
\end{mathpar}
\caption{Bidirectional term typing rules}
\label{twoltt-typing}
\end{figure}

\subsection{Computation rules and models of dependent
type theories}

Note that discussion above was ambient to the choice of
formalism to present models. Now it is finally time to
fix the formalism since we're going to need it to
define the type of models $\mathcal{M}$.

In line with current state of the art, we use
Categories-with-Families \cite{castellan2020cwf} (or
CwFs for short) as it is currently one of the most
convenient frameworks to specify models of dependent
type theories. Here is a self-contatined definition of
a category with families, given for clear reference in
subsequent internalization:

\begin{definition}[Category]
A category is the following datum\footnote{Of course,
cartesian closed category (\ref{ccc}) is also a
category. We repeat parts of definitions for
self-containment reasons.}:
\begin{enumerate}
    \item A class of \textit{objects} $\mathrm{Ob}$ and
        a family of \textit{morphisms}
        $\mathrm{Hom}(S, T)$ indexed by
        $S,T\in\mathrm{Ob}$, where $f\in\mathrm{Hom}(S,
        T)$ is denoted by $f:S\to T$. In addition, $S$
        is called \textit{source} of $f$, while $T$ is
        \textit{target} of $f$.
    \item For each $A\in\mathrm{Ob}$ there's an
        \textit{identity morphism} $\mathrm{id}_A:
        A\to A$; for each $f:A\to B$ and $g:B\to C$,
        there's their \textit{composition} $g\circ
        f:A\to C$. If a source-target of an identity
        morphism is clear from the context, the
        subscript is elided. Also, the following
        properties hold:
        \begin{align*}
            \mathrm{id}\circ f = f = f\circ\mathrm{id}
            & &
            (f \circ g) \circ h = f \circ (g \circ h)
        \end{align*}
        whenever all compositions are well-defined.
\end{enumerate}
\end{definition}

\begin{definition}[Terminal object]
    Terminal object of category $\mathcal{C}$ is an
    object $\top\in\mathrm{Ob}_\mathcal{C}$ such that
    for each $A\in\mathrm{Ob}_\mathcal{C}$ there's a
    \textbf{unique} morphism $!_A:A\to\top$.
\end{definition}

\begin{definition}[Presheaf]
    A presheaf $P$ on a category $\mathcal{C}$ (written
    as $P:\mathcal{C}^\mathrm{op}\to\mathrm{Set}$) is a
    mapping which takes objects of $\mathcal{C}$ to
    sets, and morphisms of $\mathcal{C}$ to functions
    between sets, so that if $f:A\to_\mathcal{C} B$,
    then $P(f)$ is a function from $P(B)$ to $P(A)$.
    Also, the following properties hold:
    \begin{align*}
        P(\mathrm{id}_A)=\mathrm{id}_{P(A)} & &
        P(g\circ f)=P(f)\circ P(g)
    \end{align*}
    Where $\mathrm{id}_{P(A)}$ is an identity function
    of a set $P(A)$ and $P(f)\circ P(g)$ is a
    composition of functions $P(f)$ and $P(g)$.
\end{definition}

\begin{definition}[Category of elements]
    Given a presheaf $P:\mathcal{C}^\mathrm{op}\to
    \mathrm{Set}$, a category of elements
    $\int_\mathcal{C}P$ is a category where objects are
    pairs $(A, x)$ such that
    $A\in\mathrm{Ob}_\mathcal{C}$ and $x\in P(A)$, and
    morphisms from $(A, x)$ to $(B, y)$ are morphisms
    $f:A\to_\mathcal{C} B$ such that $P(f)(y)=x$.
\end{definition}

\begin{definition}[Category with families]
    \label{cwfdef}
    A category with families is the following datum:
    \begin{enumerate}
        \item A \textit{category of contexts}
            $\mathcal{C}$ which is a category with a
            terminal object;
        \item A \textit{presheaf of types}
            $\mathrm{Ty}:\mathcal{C}^\mathrm{op}\to
            \mathrm{Set}$. Given a morphism $\gamma:
            \Gamma\to_\mathcal{C}\Delta$ and a
            \textit{type} $T\in\mathrm{Ty}(\Delta)$,
            element $\mathrm{Ty}(\gamma)(T)\in
            \mathrm{Ty}(\Gamma)$ is usually written as
            $T[\gamma]$ for brevity.
        \item A \textit{presheaf of terms}
            $\mathrm{Tm}:\left(\int_\mathcal{C}
            \mathrm{Ty}\right)^\mathrm{op}\to
            \mathrm{Set}$. Same as with types,
            $\mathrm{Tm}(\gamma)(t)\in\mathrm{Tm}(
            \Gamma, T)$ is usually written as
            $t[\gamma]$ for brevity.
        \item For each
            $\Gamma\in\mathrm{Ob}_\mathcal{C}$ and
            $T\in\mathrm{Ty}(\Gamma)$, there is an
            \textit{extended context} $\Gamma.T\in
            \mathrm{Ob}_\mathcal{C}$ together with a
            \textit{weakening morphism} $p:\Gamma.T\to
            \Gamma$ and \textit{variable element}
            $q:\mathrm{Tm}(\Gamma.T, T[p])$. In
            addition, for all $\Delta$, $\gamma:\Delta
            \to\Gamma$ and $t\in\mathrm{Tm}(\Delta,
            T[\gamma])$, there is a unique
            \textit{extended substitution} $\langle
            \gamma,t\rangle:\Delta\to\Gamma.T$ such
            that the following holds:
            \begin{align*}
                p\circ\langle\gamma,t\rangle=\gamma & &
                q[\langle \gamma, t \rangle] = t
            \end{align*}
\end{enumerate}
\end{definition}

A definition of a Cwf (def. \ref{cwfdef}) clearly
captures the bare-bones structure of a dependently
typed theory: objects of $\mathcal{C}$ represent
contexts, while morphisms represent substitutions
mapping one context to another. Substitution also acts
both on types and terms; given a correct type $T$ in
context $\Gamma$, one can extend a context to gain new
context $\Gamma,x:T$ in which there's a well-typed term
$x:T$; there's a way to obtain substitutions into an
extended context by providing a term suitable for $x$.

This base definition can be readily extended to
incorporate type (and term) formers of any system. A
CwF for $\twoLTT$, then, has to include $\mathrm{Ty}$-
and $\mathrm{Tm}$- formers for analogues of universes,
$\Pi$- and $\Sigma$- types, natural numbers and
identity types. Interpretation for $\Box_\ell$ is not
needed since we succesfully prevented boxed types from
appearing inside boxed terms via rigid \faLock-ing
discipline\footnote{That's also the reason we dubbed
this system \textit{almost} reflective, since its
internal reinterpretation doesn't cover modalities and
reinterpretation itself.}. Thus, contexts which need to
be interpreted are constructed via vanilla context
extensions from empty context, and a usual CwF is
enough for our purposes. Since actual internalized
definition of our CwF largely repeats typing and
computation rules from
Figures~\ref{twoltt-wf}-\ref{twoltt-typing} and
Figures~\ref{twoltt-typcomp}-\ref{twoltt-comp} in
internal language, we do not provide it here; to get a
bit more details, you can refer to our
Appendix~\ref{twoltt-internal-model-def}.

Finally, on to computation rules as given in
Figures~\ref{twoltt-typcomp}-\ref{twoltt-comp}.
Mostly they are standard MLTT computation rules for
bidirectional typing (without $\eta$-equality for now).
Importantly enough, we \textbf{do not compute inside
locked terms and inside boxed types} in order to not
stumble into a variant of our no-go theorem \ref{no-go}
again (modal variant is spelled out in discussion, see
Section~\ref{global-elements}).

\begin{figure}
\centering
\begin{mathpar}
    \inferrule{t\leadsto t'}
    {\mathrm{El}(t)\leadsto\mathrm{El}(t')}
    \and
    \inferrule{}{\mathrm{El}(\mathtt{U}_\kappa)\leadsto
        \mathcal{U}_\kappa}
    \and
    \inferrule{}{\mathrm{El}(\mathtt{Nat}_\kappa)
        \leadsto\mathrm{Nat}_\kappa}
    \and
    \inferrule{}{\mathrm{El}(\mathtt{Sigma}(x:s)t)
        \leadsto(x:\mathrm{El}(s))\times\mathrm{El}(t)}
    \and
    \inferrule{}{\mathrm{El}(\mathtt{Pi}(x:s)t)\leadsto
            (x:\mathrm{El}(s))\to\mathrm{El}(t)}
    \and
    \inferrule{}{\mathrm{El}(\mathtt{Id}\,t\,t_1\,t_2)  
        \leadsto t_1=_{\mathrm{El}(t)}t_2}
    \and
    \inferrule{}{\mathrm{El}(\mathtt{Box}_\kappa t)
        \leadsto\Box_\kappa\mathrm{El}(t)}
    \and
    \inferrule{S\leadsto S'}
        {(x:S)\times T\leadsto (x:S')\times T}
    \and
    \inferrule{T\leadsto T'}
    {(x:S)\times T\leadsto(x:S)\times T'}
    \and
    \inferrule{S\leadsto S'}
    {(x:S)\to T\leadsto (x:S')\to T}
    \and
    \inferrule{T\leadsto T'}
    {(x:S)\to T\leadsto(x:S)\to T'}
    \and
    \inferrule{T\leadsto T'}
        {t_1=_T t_2\leadsto t_1=_{T'} t_2}
    \and
    \inferrule{t_1\leadsto t_1'}
        {t_1=_T t_2\leadsto t_1'=_T t_2}
    \and
    \inferrule{t_2\leadsto t_2'}
        {t_1=_T t_2\leadsto t_1=_T t_2'}
    \and
    \inferrule{}{
        t\llbracket T\rrbracket\leadsto
        t.\mathrm{Tm}(t.\cdot, t\llbracket\cdot\vdash
        T\rrbracket)}
\end{mathpar}
    \caption{$\twoLTT$ typal computation rules}
\label{twoltt-typcomp}
\end{figure}

\begin{figure}
\centering
\begin{mathpar}
    \inferrule{t\leadsto t'}{S\,t\leadsto S\,t'}
    \and
    \inferrule{t\leadsto t'}{
        \mathrm{natind}_{x.T}t\{s_0,s_S\}\leadsto
        \mathrm{natind}_{x.T}t'\{s_0,s_S\}
    }
    \and
    \inferrule{}{\mathrm{natind}_{x.T}
        (0:\mathrm{Nat}_\ell)\{s_0,s_S\}
        \leadsto (s_0:T[(0:\mathrm{Nat}_\ell)/x]
        )}
    \and
    \inferrule{}{\mathrm{natind}_{x.T}
        (S\,t:\mathrm{Nat}_\ell)\{s_0,s_S\}
        \leadsto (s_S:(x:\mathrm{Nat}_\ell)\to T\to
            T[(S x)/x])\,t\,(\mathrm{natind
            }_{x.T}(t:\mathrm{Nat}_\ell)\{s_0,s_S\})}
    \and
    \inferrule{s\leadsto s'}{\mathtt{Sigma}(x:s)t
        \leadsto\mathtt{Sigma}(x:s')t}
    \and
    \inferrule
        {t\leadsto t'}
        {\mathtt{Sigma}(x:s)t\leadsto
            \mathtt{Sigma}(x:s)t'}
    \and
    \inferrule{s\leadsto s'}{(s,t)\leadsto(s',t)}
    \and
    \inferrule{t\leadsto t'}{(s,t)\leadsto(s,t')}
    \and
    \inferrule{t\leadsto t'}{\pi_i t\leadsto \pi_i t'}
    \and
    \inferrule{}
    {\pi_i((t_1,t_2):(x:S)\times T)\leadsto t_i}
    \and
    \inferrule{s\leadsto s'}{\mathtt{Pi}(x:s)t
        \leadsto\mathtt{Pi}(x:s')t}
    \and
    \inferrule
        {t\leadsto t'}
        {\mathtt{Pi}(x:s)t\leadsto
            \mathtt{Pi}(x:s)t'}
    \and
    \inferrule
        {t\leadsto t'}
        {\lambda x.t\leadsto\lambda x.t'}
    \and
    \inferrule{t\leadsto t'}{t\,s\leadsto t'\,s}
    \and
    \inferrule{s\leadsto s'}{t\,s\leadsto t\,s'}
    \and
    \inferrule{}
        {((\lambda x.t):(x:S)\to T)s\leadsto t[(s:S)/x]}
    \and
    \inferrule{t\leadsto t'}{\mathtt{Id}\,t\,t_1\,t_2
        \leadsto\mathtt{Id}\,t'\,t_1\,t_2}
    \and
    \inferrule{t_1\leadsto t_1'}{
        \mathtt{Id}\,t\,t_1\,t_2\leadsto
        \mathtt{Id}\,t\,t_1'\,t_2
    }
    \and
    \inferrule{t_2\leadsto t_2'}{
        \mathtt{Id}\,t\,t_1\,t_2\leadsto
        \mathtt{Id}\,t\,t_1\,t_2'
    }
    \and
    \inferrule{t\leadsto t'}{\mathrm{J}_{xy.S}\,t\,s
        \leadsto\mathrm{J}_{xy.S}\,t'\,s}
    \and
    \inferrule{}{\mathrm{J}_{xy.S}(\mathrm{refl}:
        t_1=_T t_2)\,s\leadsto(s:S[t_2/x,
            (\mathrm{refl}:t_1=_T t_2)/y])}
    \and
    \inferrule{t\leadsto t'}{(t:T)\leadsto(t':T)}
    \and
    \inferrule{t\leadsto t'}
    {[t]_\text{\faUnlock}\leadsto[t']_\text{\faUnlock}}
    \and
    \inferrule{}{[[t]_\text{\faLock}:\Box_\kappa T
    ]_\text{\faUnlock}\leadsto t:T}
    \and
    \inferrule{}{[[t]_\text{\faUnlock}]_\text{\faLock}
    \leadsto t}
    \and
    \inferrule{s\leadsto s'}{s\llbracket t\rrbracket
    \leadsto s'\llbracket t\rrbracket}
    \and
    \inferrule{t\leadsto t'}{s\llbracket t\rrbracket
    \leadsto s\llbracket t'\rrbracket}
    \and
    \inferrule{}{
        s\llbracket[t]_\text{\faLock}:\Box_\kappa T
        \rrbracket\leadsto
        s\llbracket\cdot\vdash t\Leftarrow T\rrbracket
    }
    \and
    \inferrule{ }{t:T\leadsto t}(t\text{ is
    synthesizing})
\end{mathpar}
\caption{$\twoLTT$ term computation rules}
\label{twoltt-comp}
\end{figure}

Also we need to give the rules for computing
reinterpretation; as it turns out,
\textit{reinterpreting terms as-is} allows to define
reinterpretation not as a redex, but rather as a couple
of total functions on types and terms, just as it was in
a non-dependently-typed case (Section~\ref{SSTLC}). To
be precise, we need two mutually recursive functions
for terms:
$m\llbracket\Gamma\vdash t\Leftarrow T\rrbracket$
computes on checkable terms, while $m\llbracket\Gamma
\vdash t\Rightarrow\rrbracket$ computes on synthesizing
terms and outputs an inferred type; both compute
elements of corresponding $\mathrm{Tm}(\Gamma,\_)$-s
from a CwF $m$. $m\llbracket\Gamma\vdash T\rrbracket$
computes an element of $\mathrm{Ty}(\Gamma)$. Another
auxiliary function is required to compute chain of
equalities to handle conversion rule; this function has
two subtle points:

\begin{itemize}
    \item At this point, we don't know if system is
        strongly normalizing yet, therefore we
        potentially might encounter an infinite chain of
        equalities.
    \item In systems where equality is proof-relevant
        (and we want to be able to develop in that
        direction), different reduction chains might
        lead to different equality witnesses, so such a
        function is potentially not well-defined.
\end{itemize}

Our solution is the following: let us fix a certain
evaluation strategy (for example, strict evaluation)
and employ it in normalization. Since terms we reduce
to find a witness are vanilla MLTT terms, they are
strongly normalizing, so witness extraction terminates.

TODO: provide interpretation definition.

\subsection{Metatheoretical properties of
\texorpdfstring{$\twoLTT$}{[[2]]LTT}}
\label{twoLTT-properties}

Once again, we only try to make sure of the most
essential properties: that computations are confluent,
that types during computations are preserved, all
computations terminate and that resulting type system
is a conservative extension of MLTT.

\begin{theorem}[Church-Rosser theorem for $\twoLTT$]
    Reflexive transitive closure of $(\leadsto)$ has
    Church-Rosser property.
\end{theorem}
\begin{proof}
    Classic proof via parallel reduction relation
    [CITE]. Since interpretation computes as a
    recursive functions on open terms, it commutes with
    substitution, which is crucial for confluence.

    TODO: write out parallel reduction?
\end{proof}

\begin{theorem}[Type preservation for $\twoLTT$]
    If $\Gamma\vdash t\Leftarrow T$ and $t\leadsto t'$,
    then $\Gamma\vdash t'\Leftarrow T$; analogously for
    $\Gamma\vdash t\Rightarrow T$.
\end{theorem}
\begin{proof}
    Note that, analogous to other dependently typed
    systems, a context reduction lemma holds which is
    crucial to handle reduction in domain of
    $\Pi$-types (where both $\Leftarrow$ and
    $\Rightarrow$ can stand in place of a colon):
    \[
        \dfrac{
            \Gamma,x:S,\Delta\vdash_\ell t:T
            \quad S\leadsto S'
        }{\Gamma,x:S',\Delta\vdash_\ell t:T}
    \]

    Then the proof proceeds by structural induction on
    terms and $(\leadsto)$. Note that, in a
    bidirectional system, one has to take care to
    structure rules in such a way that substitutions
    substitute terms that are inferred, so that
    inferrable terms stay inferrable and checkable
    terms stay checkable.

    An important difference from other dependently
    typed systems is, of course, interpretation. For
    this case, analogous to the proof of type
    preservation for $\SSTLC$
    (Section~\ref{SSTLC-properties}), a vital lemma
    holds:
    \[  \dfrac{
            \Gamma.{}_\ell\text{\faLock},\Delta
                \vdash_\kappa t\Leftarrow T
            \quad\Gamma\vdash_\ell s\Rightarrow
                \mathcal{M}_{\ell_0,\ldots,\ell_\kappa
                }^{\ell'_0,\ldots,\ell'_\kappa}
        }{\Gamma\vdash_\ell s\llbracket\Delta\vdash t
        \Leftarrow T\rrbracket\Rightarrow
            s.\mathrm{Tm}(s\llbracket\Delta\rrbracket,
                s\llbracket\Delta\vdash T\rrbracket)}
    \]
    Along with similar lemmas for
    $\_\llbracket\_\vdash\_\Rightarrow\rrbracket$ and
    $\_\llbracket\_\vdash\_\rrbracket$.
\end{proof}

\begin{theorem}[Identity model]
    For each $\ell$, there is an identity model
    $I_\ell:\mathcal{M}_{0,\ldots,\ell}^{1,\ldots,
    \ell+1}$ with a property that $I_\ell\llbracket
    [t]_\text{\faLock}:\Box_\kappa T\rrbracket\equiv
    t:T$ for each closed $t$.
\end{theorem}
\begin{proof}
    Take a universe at appropriate level as type of
    contexts; $\mathrm{Ty}_\ell(\Gamma)=
    \Gamma\to\mathcal{U}_\ell$;
    $\mathrm{Tm}_\ell(\Gamma, T)=(\gamma:\Gamma)\to
    T\,\gamma$. Define all of the remaining term
    interpretations to reinterpret to the same terms.
\end{proof}

\begin{theorem}[Strong normalization for $\twoLTT$]
    Closed terms typable in $\twoLTT$ are strongly
    normalizing with respect to $(\leadsto)$.
\end{theorem}
\begin{proof}
    We could reprove normalization from scratch, just
    how we did with $\SSTLC$. However, here we can
    provide a much simpler proof which delegates
    normalization to another type system. To be more
    precise, this system $\mathcal{S}$ is Martin-L\"of
    type theory with $\omega+\omega$ universes and
    first-class universe polymorphism allowed to range
    over first $\omega$ universes; we know it to be
    strongly normalizing [CITE]. Then, we define a
    translation function $[\_]$ which transforms types
    and terms of $\twoLTT$ into types and terms of
    $\mathcal{S}$: it preserves all vanilla grammar
    formers and only translates lock
    $[\_]_\text{\faLock}$ into a function from the type
    of models $\mathcal{M}_{\ell_0,\ldots,\ell_\kappa
    }^{\ell'_0,\ldots,\ell'_\kappa}$ into interpreted
    type. Since $\kappa$ is fixed, first-class universe
    polymorphism is enough to translate it into a
    function which accepts models for \emph{any} list
    of levels of fitting length. Then, unlocking
    $[\_]_\text{\faUnlock}$ translates into application
    of this function to identity model of fitting
    level; interpretation translates into application.

    Note that $t\leadsto t'$ implies
    $[t]\leadsto^+[t']$; this concludes the proof since
    any infinite chain of reductions in our system
    would imply an existence of infinite chain of
    reductions in $\mathcal{S}$.
\end{proof}

\begin{theorem}[Conservativity over MLTT]
    For MLTT context $\Gamma$, MLTT type $T$ and
    $\twoLTT$ term $t$, if
    $\Gamma\vdash t\Leftarrow T$, then
    $\Gamma\vdash_\mathrm{MLTT}\mathrm{nf}(t)\Leftarrow
    T$.
\end{theorem}
\begin{proof}
    Thanks to type preservation, normal form
    $\mathrm{nf}(t)$ is typable in $\twoLTT$. Since
    $\Gamma$ does not contain neither $\Box$-es nor
    \faLock-s, any subterm of type $\Box_\kappa T$
    has form $[\_]_\text{\faLock}$. Since $T$ is an MLTT
    type, every such locking is wrapped either in a
    $[\_]_\text{\faUnlock}$ or in a $\llbracket\_
    \rrbracket$. Second case is a redex and is
    impossible in a normal form; $[\_]_\text{\faUnlock}$
    is possible only in a locked context, so that means
    that there's another $[\_]_\text{\faLock}$ above it.
    Therefore $\mathrm{nf}(t)$ has no subterms of type
    $\Box_\kappa T$; then it has no unlocking and
    interpretation subterms and is an MLTT term and is
    typable in bidirectional MLTT (it is obtained from
    $\twoLTT$ by getting rid of rules involving
    modalities, locking and reinterpretation).
\end{proof}

\subsection{Useful reinterpretations inside
\texorpdfstring{$\twoLTT$}{[[2]]LTT}}
\label{twoLTT-applications}

Here we list some of the language features which
reinterpretation enables without any additional language
support.

\subsubsection{Crude, but effective stratification}
\label{twoLTT-displacement}

Due to McBride \cite{McBride:Crude2011,
McBride:Crude2012}, a technique of crude, but effective
stratification works as follows: instead of providing
complete support for level polymorphism, let's consider
a \textit{displacement operator} which, given a term $t$
using universes $U_i,U_j,U_k,\ldots$, produces a term
which is equal to $t$ except for universes which are
replaced with $U_{i+1},U_{j+1},U_{k+1},\ldots$. The
displaced term is also well-typed, however it is defined
<<one level higher>>. While not directly definable as a
single operator inside $\twoLTT$, for each
$\kappa,\ell$, one can define displacement operators
$\uparrow_\ell^\kappa:\mathcal{M}_{\kappa+0,
\kappa+1,\ldots,\kappa+\ell}^{\kappa+1,\kappa+2,\ldots,
\kappa+\ell+1}$ that shifts a boxed definition which
uses universes up to $\mathcal{U}_\ell$ by $\kappa$
levels up. In this model, category of contexts basic
presheaves are shifted by $\kappa$ in comparison to an
identity model; all other term interpretations are
identical to the ones given in an identity model.

A bit more details on definition of
$\uparrow_\ell^\kappa$ are given in
Appendix~\ref{twoltt-displacement-model-def}.

\subsubsection{Internal external parametricity}
\label{twoLTT-parametricity}

Demonstrated by Reynolds \cite{reynolds1983types},
parametricity is a property of a type system that, in
simple words, <<polymorphic function cannot
meaningfully depend on generic inputs>>. Thus any
polymorphic function of type $\sigma$ satisfies a
certain <<free theorem>> $\llbracket \sigma
\rrbracket$ defined via relational interpretation of a
type.

In recent works \cite{Bernardy:Parametricity2010,
Altenkirch:Parametricity2024}, a streamlined definition
is given for parametricity interpretation of
dependently-typed systems dubbed <<external
parametricity>>. Once again, this relational
interpretation is definable in our system
\textit{internally} with the help of a relational model
$\mathrm{Para}_\ell:\mathcal{M}_{1,\ldots,\ell+1}^{2,
\ldots,\ell+2}$. The shift by one level happens because
type $T:\mathcal{U}_\ell$ in such a model is
interpreted by a triple $(T_1\,T_2:\mathcal{U}_\ell)
\times(T_1\to T_2\to\mathcal{U}_\ell)$. Since
interpreted functions do not depend on local context,
we sidestep issues with relational context highlighted
by Altenkirch and others
\cite{Altenkirch:Parametricity2024}, thus result of the
interpretation is a proper proof of a free theorem.

A bit more details on definition of
$\mathrm{Para}_\ell$ are given in
Appendix~\ref{twoltt-parametricity-model-def}.

\subsubsection{Synthetic mathematics}
\label{twoLTT-synthetic}

As thoroughly discussed in a recent paper by Nawrocki
and others \cite{Nawrocki:Lean2026}, thanks to MLTT
being an internal language of certain
$(\infty,1)$-categories, we can do \textit{synthetic
mathematics} inside it: prove a theorem inside MLTT
about types; interpret it in a model of MLTT and yield
a proof about objects in that model. It is easy to see
that reinterpretation operator we have in $\twoLTT$ is
sufficient to do synthetic mathematics in it and
immediately apply results interpreted in a suitable
model.
